\uuid{bwA8}
\exo7id{3406}
\titre{exo7 3406}
\auteur{quercia}
\organisation{exo7}
\datecreate{2010-03-10}
\isIndication{false}
\isCorrection{false}
\chapitre{Matrice}
\sousChapitre{Matrice et application linéaire}
\cpge{mpsi-mp2i.matrices-applications-lineaires.a, pcsi.matrices-applications-lineaires.a, ptsi.matrices-applications-lineaires.a, mpsi-mp2i.changements-bases.a, pcsi.matrices-applications-lineaires.c, ptsi.matrices-applications-lineaires.c, mpsi-mp2i.applications-lineaires.b, pcsi.applications-lineaires.b, ptsi.applications-lineaires.b, bcpst1.matrices}
\module{Algèbre}
\niveau{L1}
\difficulte{}

\contenu{
\question{
Soit $A \in \mathcal{M}{n+1}(\Q)$ telle que $a_{ij} = C_{j-1}^{i-1}$.
Interpréter $A$ comme la matrice d'un endomorphisme simple de $\Q_n[X]$.
En déduire la matrice $A^{-1}$.
}
}
